% !TEX root = ../../main.tex

\section{连续识别方案探索与局限性}

除采用“词级分类模型 + 滑动窗口 + 词段后处理”方案外，本项目在实验过程中也对更直接的连续手语识别方案进行了探索。
连续识别方案的基本思想是让模型直接从一段较长的视频特征序列中输出 Gloss 序列，而不是先将视频切分为大量候选窗口，再通过分类和后处理得到词段。
这类方案通常需要使用 CTC 等序列建模方法，以处理输入帧序列和输出 Gloss 序列长度不一致的问题。

在探索阶段，项目曾基于 CE-CSL 数据进行 CTC 方向实验。
一类实验采用 Transformer Encoder + CTC 结构，输入特征使用 raw\_delta 模式，并通过位置编码、padding mask 和 CTC 解码完成序列预测。
该方案能够完成训练并输出序列，但实验结果并不理想，主要问题表现为预测序列偏短、删除错误较多，说明模型虽然具备一定学习能力，但在当前特征、数据规模和训练配置下，仍难以稳定输出完整 Gloss 序列。
因此，该方向没有被纳入 HearBridge 当前系统主线。

另一类探索使用 Conformer Encoder + CTC 结构，并在 CE-CSL 数据上引入 controlled vocabulary，将低频 Gloss 映射为 \texttt{<unk>}，以降低完整词表带来的稀疏性问题。
该方案在训练脚本中记录 train\_loss、dev\_loss 和 dev\_TER，并保存最佳 dev\_TER、最佳 dev\_loss 和最后一轮 checkpoint，适合作为后续连续识别研究的实验基础。
但由于该实验依赖较大规模连续手语数据、训练成本较高，且与 HearBridge 当前 25 词短句识别任务和系统接入链路并不完全一致，因此本项目最终没有将其作为毕业设计系统的主要实现方案。

相比之下，滑动窗口词级识别方案更适合本项目阶段目标。
一方面，WLASL 小词表能够较方便地构建 $20\times232$ 词级训练样本，模型训练和调试成本较低；
另一方面，滑动窗口方案可以在没有精确句子级边界标注的情况下，对短视频中的局部词段进行密集识别，并通过 NMS 和 DeepSeek 候选排序降低插入误报和语义不自然问题。

当前方案仍然存在明显局限。
首先，滑动窗口方式依赖固定窗口长度和步长，动作边界仍然通过后处理间接确定，无法像完整连续识别模型那样从序列层面统一建模。
其次，词级分类模型的识别范围受限于小词表，无法覆盖开放场景下的大规模 Gloss。
再次，DeepSeek 语义修正只能在程序生成的候选集合内选择更合理的序列，不能弥补视觉模型完全未识别出的关键词。
